\chapter{Theoretical Background}\label{ch:background}

\section{The ARQMath collection and tasks}\label{sec:bg-arqmath}

Every experiment in this thesis is run on the ARQMath test collection, so its
properties determine what can and cannot be measured here. This section describes
the corpus, the two retrieval tasks it defines, and the relevance criteria behind
each.

ARQMath is a community question answering (CQA) collection built from a March
2020 snapshot of Mathematics Stack Exchange, covering questions and answers
posted between 2010 and 2018~\cite{zanibbi2020arqmath}. It contains roughly
one million questions and 28 million formula instances, each formula carrying
a unique identifier and distributed alongside pre-computed \LaTeX{},
Presentation MathML and Content MathML representations. All three lab editions
--- ARQMath-1 (2020), ARQMath-2 (2021) and ARQMath-3 (2022) --- search the same
collection of posts, and differ in their topics, their relevance judgments, and the
corrections the organisers made to the formula
representations~\cite{mansouri2022arqmath3}.

Answer retrieval (ARQMath Task~1) asks a system to rank up to 1{,}000 answer
posts for a mathematical question. Each edition draws its topics from questions
posted to Mathematics Stack Exchange after the corpus snapshot, so their own
answers are not in the collection --- ARQMath-3's come from 2021 --- and only
answer posts from 2010--2018 are searchable. Candidate topics had to contain at
least one formula and have at least one moderator-annotated duplicate question
inside the collection, so that assessment effort was not spent on topics with no
relevant answer at all~\cite{mansouri2022arqmath3}. For ARQMath-3, 100 topics
were released, 80 were assessed, and two were removed because only a single
answer had been judged relevant, leaving 78 topics in the evaluated test
collection~\cite{mansouri2022arqmath3}.

Formula search (ARQMath Task~2) takes a single formula from one of those
questions as the query and returns formula instances drawn from both questions
and answers. Both post types are searched because a retrieved formula is only
ever a route to text: one occurring in an answer may sit in a worked solution,
and one occurring in another question routes the searcher to that question's own
answers, so restricting the search to answers would discard half the routes to
useful material. The ARQMath-3 test collection has 76 assessed
topics~\cite{mansouri2022arqmath3}. Relevance is contextual: a candidate formula
is judged inside the post it appears in, and the query formula inside the
question it was taken from. Evaluation is performed over \emph{visually
distinct} formulae rather than instances --- those that render identically are
clustered under a shared visual identifier and credited once --- so retrieving
twenty copies of the same expression earns no additional
reward~\cite{mansouri2021effects}. A third pilot task, open-domain question
answering, was introduced in 2022 and is not used here.

Both tasks use the same four relevance levels but define them differently, and
the difference is the point. In answer retrieval, an answer is High if it is
sufficient to answer the complete question on its own, Medium if it provides
some path towards a solution, Low if it only helps find or interpret an answer,
and otherwise Not relevant --- a level the criteria explicitly extend to any
post that restates the question without adding
information~\cite{mansouri2022arqmath3}. In formula search, a formula is High if
it is as good as an exact match to the query would have been, Medium if useful
but less so, and Low if there is some chance of finding something useful through
it~\cite{mansouri2022arqmath3}. Answer-retrieval relevance is thus the extent to
which an answer resolves the question, and formula search asks how likely a
formula is to lead a searcher to text that helps answer the question its query
was taken from. Neither is structural similarity: a formula matching the query
symbol-for-symbol can be judged non-relevant if its post is unhelpful, and a
structurally dissimilar formula can be highly relevant. They are not equally far
from it, though. The formula-search scale uses an exact match as its reference
point: a formula earns the top grade when finding it would have been as useful
as finding the query formula itself. A system that ranks by structural closeness
is ranking candidates by how near they come to that exact match, so the system
and the scale are pointed at the same thing. The answer-retrieval scale never
mentions formulas at all.

Two binarisations of these levels recur throughout. \emph{H+M} counts only High and
Medium as relevant, matching the usual TREC notion, while \emph{H+M+L} additionally
counts Low; binary measures in this thesis use H+M, that is, relevance level 2
(Section~\ref{sec:bg-retrieval-metrics}).

Judgments were produced by pooling over the participants' \emph{runs}. A run is
the ranked list of answers a participating system submits for every topic; each
team designates one of its runs as primary and the others as alternate. For each
topic the organisers took the top 45 results of every primary run and the top 20
of every alternate run, merged and deduplicated them, and had trained assessors
judge the result, yielding pools averaging 464 answers per topic for answer
retrieval in ARQMath-3~\cite{mansouri2022arqmath3}. The judgments are
distributed as \emph{qrels} (query relevance judgments), one entry per judged
topic--document pair carrying its grade, so a qrels file records the pool as
well as the grades: a document is judged exactly when it has an entry. The pool
is therefore defined by what the 2022 participants retrieved, and any document
outside it is unjudged rather than non-relevant. Such documents are called
\emph{holes}, and the fraction of a ranking's top-ranked results that consists
of them is its \emph{hole rate}, made precise at a cutoff in
Section~\ref{sec:bg-retrieval-metrics}; incomplete judgments of this kind are a
long-standing problem in pooled evaluation~\cite{buckley2004incomplete}. The
organisers' response was to adopt measures that ignore holes rather than scoring
them as irrelevant: the primary measure for both tasks is $\mathrm{nDCG}'$, the
prime variant of Sakai and Kando~\cite{sakai2008prime}, with $\mathrm{MAP}'$ and
$\mathrm{P}'@10$ as secondary measures. That convention suffices for the systems
that were pooled, but not for a system built afterwards, which retrieves
documents no 2022 participant retrieved and therefore carries a far higher hole
rate; Sections~\ref{sec:bg-retrieval-metrics} and~\ref{sec:met-judge-validation}
take up that gap.

\section{Formula representations}\label{sec:bg-formula-representations}

A formula written in \LaTeX{} is a string, and string matching cannot see that
$x - y^2$ and $a - b^2$ are the same expression under different names. Retrieving
on structure therefore requires parsing the formula into a tree first. \LaTeX{} is
presentation source: it fixes which symbols appear and where they sit on the page,
but encodes no structure a parser can read off directly. \emph{Mathematical
markup} does encode it, and the standard for it is MathML, an XML language for
mathematics that comes in two flavours. Those two flavours support two different
trees, and the difference between them runs through the rest of this thesis:
Presentation MathML records how a formula is laid out on the page, while Content
MathML records what it computes. ARQMath
distributes both for every formula in the collection
(Section~\ref{sec:bg-arqmath}), so the trees can be read from the collection
rather than parsed out of \LaTeX{}.

The \emph{Symbol Layout Tree} (SLT) is derived from Presentation MathML and
encodes spatial arrangement. Its nodes are the symbols of the formula and its
edges are the spatial relations between them --- one symbol follows another on the
same baseline, or sits above it as a superscript, or below it, or within a
structure such as a radical.

The \emph{Operator Tree} (OPT) is derived from Content MathML and encodes
operator--operand structure instead. Its root is the top-level operator or
relation, its internal nodes are the remaining operators, and its leaves are
operands; each edge records which argument of its parent the child is. The OPT
recognises $y^2$ as an exponentiation however it was typeset, but no longer knows
what the expression looked like.

Both trees are linearised the same way. Following
Tangent-CFT~\cite{mansouri2019tangentcft}, which extends the symbol pairs of
Tangent-S~\cite{davila2017tangents} with the path from the root, a depth-first
traversal converts a tree into a list of \emph{path tuples}, each pairing a node
with one of its descendants, and every tuple, in either representation, has the
same four positions:

\begin{description}[leftmargin=4.2em,style=sameline,itemsep=0.15em,topsep=0.4em]
  \item[\textmd{\texttt{[0]}}] a node;
  \item[\textmd{\texttt{[1]}}] a node that descends from it;
  \item[\textmd{\texttt{[2]}}] the edge labels on the path from the first to the second;
  \item[\textmd{\texttt{[3]}}] the path from the root of the tree to the first node.
\end{description}

\noindent
What fills those positions is what separates the two representations. In an SLT,
positions \texttt{[0]} and \texttt{[1]} hold a symbol written with its type ---
\texttt{V!x} for the variable $x$ and \texttt{N!2} for the number $2$, with
\texttt{!} separating type from value --- while operators such as the minus sign
appear as the bare symbol \texttt{-}. Position \texttt{[2]} names the spatial
relation: \texttt{n} for the next symbol on the same baseline, \texttt{a} for
above, \texttt{b} for below, \texttt{w} for within, among others. Position
\texttt{[3]} is then the sequence of those relations taken from the root, so
\texttt{nn} is two steps along a baseline, and a tuple whose first node \emph{is}
the root carries \texttt{-}. In $x - y^2 = 0$, for instance, the tuple
\texttt{(V!y, N!2, a, nn)} records that the $2$ sits above the $y$, and that the
$y$ lies two steps along the baseline from the root.

An OPT fills the same four positions differently. Its operators are marked for
whether the order of their arguments matters: the subtraction in $x - y^2 = 0$
becomes \texttt{O!minus}, because $x - y^2$ and $y^2 - x$ are different, and the
equality becomes \texttt{U!eq}, because $x - y^2 = 0$ and $0 = x - y^2$ are not.
Position \texttt{[2]} gives the argument index, \texttt{0} for the first argument
and \texttt{1} for the second, and position \texttt{[3]} the sequence of those indices
from the root, so \texttt{01} is the first argument of the root followed by the
second of its child. In the same formula, \texttt{(O!SUP, N!2, 1, 01)} records that
the $2$ is the second argument of the exponentiation, which itself sits at
\texttt{01}: it is the second argument of the subtraction, and the subtraction is
the first argument of the equality.

Both close their paths with a dummy marker, written \texttt{0!} --- at the end of
each baseline in an SLT, and at each leaf in an OPT --- so that short expressions
still produce matchable tuples.

How far from a node the traversal looks is set by a \emph{window}. With a window of
one, a node is paired only with its children, so the SLT of $x - y$ yields, besides
its end marker, \texttt{(V!x, -, n, -)} and \texttt{(-, V!y, n, n)}. With a window
of two, each node is also paired with the descendants two edges away, and $x - y$
yields \texttt{(V!x, V!y, nn, -)} as well, with position \texttt{[2]} recording
both steps. A wider window captures relations between symbols further apart, at the
cost of more tuples per formula.

The third representation, \emph{SLT-TYPE}, is the SLT with its values erased and
only the node types kept: every variable collapses to \texttt{V!} and every number
to \texttt{N!}, so that \texttt{(V!x, -, n, -)} becomes \texttt{(V!, -, n, -)}.
Operators carry no type marker in an SLT and so survive intact: what is forgotten
is which symbols the operators act on, not the operators themselves. This makes
$x - y^2$ match
$a - b^2$ exactly, which is the intent, and also makes $x + x$ indistinguishable
from $x + y$, which is the cost.

The three representations therefore sit on an abstraction ladder. SLT is the most
literal and matches on appearance; OPT ignores notation but keeps operator
identity; SLT-TYPE keeps the operators and forgets the symbols they act on. Each
representation finds more formulas that are written differently but are related to
some extent. Which step is the better trade-off cannot be settled by argument; it
has to be measured. Tangent-CFT sidesteps the choice altogether by summing the
vectors of all three representations into one~\cite{mansouri2019tangentcft}. This
thesis does the opposite: it builds a separate index for each representation and
compares them against one another, so
that the contribution of each can be seen on its own (RQ1.3).
Section~\ref{sec:bg-indexing} describes how a tuple sequence becomes a vector that
can be indexed and searched.

\section{Indexing and retrieval}\label{sec:bg-indexing}

Retrieval methods differ in how a document is represented and how its similarity
to a query is computed, and each representation needs a different kind of index to
be searched at scale. Three families bear on this thesis: \emph{sparse retrieval}
over the terms a document contains, \emph{dense retrieval} over learned vectors,
and retrieval over the structure of the formulas a document holds. All three do
the same thing --- score every document in the collection against a query and
return a ranked list --- and this section describes how each represents a
document, how each is searched, and what each can and cannot match. The third is
the only one of the three that reads a formula as anything but a string.

\emph{Sparse retrieval} scores documents by the query terms they contain, and
BM25 is the standard weighting of this kind. BM25 weights a term by how often it
occurs in the document, how rare it is across the collection, and how long the
document is; two parameters shape the weighting, $k_1$ for how quickly repeated
occurrences of a term stop adding score, and $b$ for how strongly a document is
penalised for its length~\cite{robertson2009bm25}. Matching is lexical, so a
document is scored only if it literally shares a term with the query. For
mathematical text this cuts both ways: a formula matches only when its \LaTeX{}
tokens appear verbatim, but when they do the match is exact and cheap.

\emph{Dense retrieval} replaces terms with meaning. A neural encoder maps the
query and each document to a single fixed-length vector, and similarity is the dot
product between them. The architecture is a \emph{bi-encoder}: query and document
are encoded independently, which is what makes the approach practical, since every
document can be encoded once offline and only the query has to be encoded at
search time~\cite{karpukhin2020dpr}. Encoders of this kind are
pre-trained on general text and then fine-tuned on large collections of
question--answer pairs, which is what adapts them to the genre of a community
question answering collection. Their strength is that they match paraphrases a
lexical match misses; their weakness is that a whole answer post, however long,
must be compressed into one vector.

Formula retrieval embeds the path tuples of
Section~\ref{sec:bg-formula-representations}. The tuple sequence of a formula is
treated as a sentence and each tuple as a word, an approach introduced by
Tangent-CFT~\cite{mansouri2019tangentcft}, and a \emph{fastText} model is trained
over the resulting corpus --- a word-embedding method that learns a vector for
each token from the tokens surrounding it, and that composes each token's vector
from its character n-grams as well as the token itself, rather than from the token
alone~\cite{bojanowski2017fasttext}. The training objective is skip-gram: each
tuple is trained to predict the tuples around it and to distinguish them from
randomly sampled ones, so tuples occurring in similar surroundings end up with
similar vectors. Because a token is assembled from its n-grams, a tuple never seen
in training still receives a vector. That matters here because mathematical
expressions are close to unique, and a query formula absent from the collection is
the normal case. A formula is represented by the mean of its tuple vectors, and
two formulas are compared by the cosine of theirs.

Dense and formula retrieval need a different index from a sparse one. An inverted
index scores only the documents sharing a term with the query, whereas a vector
representation gives every document a non-zero similarity, making retrieval a
nearest-neighbour search over the whole collection. Exhaustive search is exact but
costs time proportional to the number of vectors, which is affordable for the
answer posts and not for the roughly 28 million formula instances. \emph{Approximate
nearest-neighbour} search buys that cost back: the vectors are grouped into cells
around learned centroids so that only the cells nearest the query are scanned, and
each vector is stored compressed so that the index fits in
memory~\cite{johnson2021faiss}. Both steps trade a little recall for a large
saving in time and space, and how much is traded is a parameter rather than a
property of the method; Section~\ref{sec:met-formula-indexing} gives the settings
used. Each of the three then returns a ranked list, but on scales that cannot be
compared with one another, which is the problem Section~\ref{sec:bg-rank-fusion}
addresses.

\section{Rank fusion}\label{sec:bg-rank-fusion}

Mapping those scores onto a common scale is the obvious repair, and it is not
free. A BM25 score is unbounded and depends on collection statistics, a dot
product depends on the norms of two embeddings, and a cosine is confined to
$[-1,1]$. Adding or averaging such scores is meaningless until they are made
commensurable, and every way of making them so is itself a modelling decision:
min--max normalisation, for instance, depends on the highest and lowest score that
happened to appear in that particular result list, so the same document can be
worth a different amount depending on what else was retrieved alongside it.
Rank-based fusion avoids the choice altogether by discarding the scores and
keeping only the position each document occupies, which is comparable across
rankings by construction.

Reciprocal Rank Fusion (RRF) is the simplest and most widely used method of this
kind. Given a set of rankings $R$, each document $d$ receives the score
\begin{equation}
\mathrm{RRF}(d) = \sum_{r \in R_d} \frac{1}{k + r(d)},
\end{equation}
where $R_d \subseteq R$ is the set of rankings in which $d$ appears, $r(d)$ is the
rank of $d$ in ranking $r$, and $k$ is a constant, set to 60 in the original
work~\cite{cormack2009rrf}. The reciprocal shape means a document gains most from
appearing near the top of a list and little from appearing deep in one, while $k$
damps the difference between the highest positions so that no single ranking can
decide the outcome by placing one document first. A document
retrieved by several rankings accumulates contributions and rises; a document
found by only one of them still enters the fused list. A ranking whose values are on a
wildly different scale cannot dominate the result, because each contribution is
bounded by $1/(k+1)$. That is precisely what makes RRF usable across sparse,
dense and formula rankings without calibrating any of them against one another.
The same peculiarity of RRF cuts the other way too. A contribution depends only on
position, so the document a ranking places first is worth $1/(k+1)$ whether that
ranking is a good one or a poor one: the rule weighs rankings equally because it
has no measure of their quality, and a ranking that retrieves badly therefore
pushes its own top documents into the fused list rather than merely failing to
improve it.

Which rankings are worth combining, and with what value of $k$, the rule does not
settle for itself; Section~\ref{sec:met-retrieval-configs} states the choices this
thesis makes. What no fusion rule settles at all is whether the list it produces
is better than the ones that went into it --- a question nothing inside a ranking
can answer, and the one Section~\ref{sec:bg-retrieval-metrics} takes up.

\section{Retrieval evaluation metrics}\label{sec:bg-retrieval-metrics}

A retrieval measure compares a ranked list of documents against the relevance
judgments for the topic that produced it. ARQMath judges on the four-level scale of
Section~\ref{sec:bg-arqmath}, so each judged document carries a grade
$g \in \{0,1,2,3\}$, and a measure turns a ranking and those grades into one number
per topic, averaged over topics to score a system.

Two kinds of measure are standard on a graded collection of this sort. One counts
how many relevant documents reach the head of a ranking; the other also reads the
order in which they arrive.

\emph{Precision at $k$} is the simpler of the two. It binarises the grades by the
H+M rule --- a document counts as relevant when $g \geq 2$ --- and reports the
fraction of the top $k$ documents that are relevant:
\begin{equation}
\mathrm{P}@k = \frac{1}{k}\,\bigl|\{d \in \text{top-}k : g(d) \geq 2\}\bigr| .
\end{equation}
It ignores both the grades above the threshold and the order within the top $k$: a
ranking with three relevant documents scores the same whether they occupy the
first three positions or the last three.

\emph{Normalised discounted cumulative gain} keeps both~\cite{jarvelin2002cumulated}.
Each document contributes a gain equal to its grade, discounted by how far down the
ranking it sits, and these are accumulated:
\begin{equation}
\mathrm{DCG}@k = \sum_{i=1}^{k} \frac{g(d_i)}{\log_2 (i+1)} ,
\end{equation}
where $d_i$ is the document at rank $i$. The discount written here is the one that
applies from rank 1~\cite{burges2005learning}, and is what \texttt{trec\_eval}
computes; the original formulation leaves the first ranks
undiscounted~\cite{jarvelin2002cumulated}. The logarithmic discount means a High
answer at rank 1 is worth more than the same answer at rank 10, which is the
property precision lacks. DCG is unbounded and depends on how many relevant
documents a topic happens to have, so it is normalised by the largest value any
ranking of that topic's judged documents could achieve --- the \emph{ideal} DCG,
obtained by sorting all judged documents for the topic by grade and taking the
first $k$:
\begin{equation}
\mathrm{nDCG}@k = \frac{\mathrm{DCG}@k}{\mathrm{IDCG}@k} .
\end{equation}
A score of 1 means the top $k$ could not have been ordered better, and 0 that
nothing relevant was retrieved. The ideal is a property of the judgments, not of the
ranking being scored: a relevant document added to a topic's judgments raises that
topic's ideal for every system at once. Two nDCG values are therefore comparable only
when they are computed against the same judgments.

\emph{Mean average precision} reads the ranking to its end rather than stopping at
a cutoff. For a single topic, average precision adds up the precision at every rank
holding a relevant document and divides by the number of relevant documents the
topic has:
\begin{equation}
\mathrm{AP} = \frac{1}{R} \sum_{i \,:\, g(d_i) \geq 2} \mathrm{P}@i ,
\end{equation}
where $R$ is that number. MAP is the mean of AP over topics. It binarises at H+M as
precision does, and rewards relevant documents that arrive early as nDCG does, but
it is the only one of the three whose value depends on the whole ranking; it is
also one of the measures ARQMath publishes for its participant
runs~\cite{mansouri2022arqmath3}.

All three measures assume that every document they read carries a grade.

\subsection*{Prime variants}

A test collection judges a pool assembled from the systems that took part in the
original evaluation. A system built afterwards retrieves documents that were never
in that pool, and those documents have no judgment at all. Treating them as
non-relevant penalises exactly the systems that retrieve differently from the pool,
which is the case relevant to this thesis (Section~\ref{sec:met-hole-problem}).

The \emph{prime} variants answer this by condensing the list: every unjudged
document is removed from the ranking, and the measure is computed over what
remains, so that the cutoff $k$ is reached using judged documents
only~\cite{sakai2008prime}. The ideal is untouched: it is still computed from every
judged document for the topic, so condensing changes what a ranking is scored on,
not what it is scored against. Written P$'$@$k$, nDCG$'$@$k$ and MAP$'$, they are the
official measures of the ARQMath labs~\cite{mansouri2022arqmath3} and are used
throughout this thesis. Sakai found condensed-list measures more robust to
incomplete judgments than the alternatives designed for the same
problem~\cite{sakai2007alternatives}, bpref among them~\cite{buckley2004incomplete}.

Condensing is not free. Removing unjudged documents implicitly assumes they would
have been graded like the judged ones around them, and where that assumption fails
the result is optimistic: condensed-list measures tend to overestimate systems that
did not contribute to the pool~\cite{froebe2023bootstrapped}. The size of the
assumption is quantified by the \emph{hole rate}, the fraction of a ranking's top
$k$ that carries no judgment,
\begin{equation}
\mathrm{hole}@k = \frac{1}{k}\,\bigl|\{d \in \text{top-}k : d \text{ unjudged}\}\bigr| ,
\end{equation}
which says how much of the ranking a prime score was in fact computed over.
Section~\ref{sec:met-judge-validation} describes the second response to the same
problem, which is to have the unjudged documents labelled rather than discarded.
Since nDCG values are comparable only under shared judgments, that response must also
settle which judgments every system is then scored against, and
Section~\ref{sec:met-evaluation-protocol} settles it.

\section{Relevance assessment and agreement}\label{sec:bg-agreement}

% Written 2026-09-19 (manifest §7a, §3.6 rows). Background only: no model, prompt, threshold or
% number of this thesis -- those are §4.4.

Section~\ref{sec:bg-retrieval-metrics} ended on a problem. A system built after a
collection was assessed retrieves answers no assessor saw, and discarding them
overestimates it. The other response is to have those answers labelled by someone
other than the original assessors, and increasingly that someone is a language model.

Whether such labels can stand in for the assessors' is a question that can be
answered empirically, and it comes in two forms that are easy to confuse. The first is
whether the two labellers give the same grade to the same answer. The second is
whether scoring systems with one labeller's grades orders them the way the other's
grades do. This section first introduces language models as relevance assessors, then gives a statistic for each of the two questions, and ends with how small a difference between two systems an evaluation can detect. The two questions have to be kept apart: a labeller can fail the first and still pass the second.

Relevance is usually judged by people, and human assessment is slow and expensive. That is why pools stop at a fixed depth
(Section~\ref{sec:bg-arqmath}). A large language model can be prompted to do the same
task: given a query and a document, assign one of the collection's relevance grades. A
model used this way is an \emph{LLM relevance judge}, or simply a \emph{judge}.

Thomas et al.\ showed that a carefully prompted model can match human relevance
preferences in web search~\cite{thomas2024llm}. UMBRELA, an open-source reproduction of
their prompt, was used to assess the TREC 2024 RAG track alongside NIST's
assessors~\cite{upadhyay2025umbrela}. Its prompt describes the task, sets out each
relevance level in turn, and forces a single grade as output. The comparison with
NIST's assessors established three things: the judge reproduced the ordering of
systems closely; it agreed with the assessors much less well on individual documents;
and it was more lenient, rating documents relevant that the assessors did not. Its
authors concluded that such a judge can replace manual judgments to measure
effectiveness ``in a coarse-grained manner'', but should not be used to separate
systems whose effectiveness is close. Each of these findings needs a statistic, and
the rest of this section supplies them.

\subsection*{Agreement on individual answers}

When two labellers grade the same items, their agreement is summarised in a
\emph{confusion matrix}. It has one row per grade the first labeller used and one
column per grade the second used, and each cell counts the items that received that
pair of grades. The diagonal holds the items on which the two agree. The share of all
items that sit on the diagonal is the \emph{observed agreement}, $p_o$.

Observed agreement flatters any labeller when one grade dominates. In a pooled
collection most judged documents are not relevant. Two labellers who never read a
document but answered ``not relevant'' every time would agree on most of them.

\emph{Chance agreement}, $p_e$, measures that floor. Call the two labellers $a$ and $b$, and write $P_a(k)$ and $P_b(k)$ for the share of items each of them assigns grade $k$. If the two assigned grades independently, at those rates, they would agree on grade $k$ with probability $P_a(k)\,P_b(k)$, and on some grade with probability
\begin{equation}
p_e = \sum_k P_a(k)\,P_b(k) .
\end{equation}
\emph{Cohen's $\kappa$} removes this floor~\cite{cohen1960kappa}:
\begin{equation}
\kappa = \frac{p_o - p_e}{1 - p_e} .
\end{equation}
The numerator is the agreement achieved beyond chance, and the denominator is the
agreement that was available beyond chance. $\kappa$ is 0 for a labeller no better
than chance, 1 for perfect agreement, and negative when the two disagree more than
chance would. On a graded scale, $\kappa$ can be computed over the grades as given,
which asks whether the two chose the same grade. It can also be computed after
binarising, for instance by the H+M rule of Section~\ref{sec:bg-arqmath}, which asks
only whether they put the item on the same side of the relevance threshold.

$\kappa$ is not a property of the labeller alone. Because $p_e$ is built from how
often each grade is used, the same observed agreement yields a smaller $\kappa$ when
one grade dominates. In that case chance alone already accounts for most of the agreement, so $p_e$ is close to $p_o$, and $\kappa$ comes out low even when the two labellers agree on most items. This is the
\emph{prevalence effect}~\cite{simwright2005kappa}, one of the two paradoxes of high
agreement but low $\kappa$~\cite{feinstein1990kappa}. Two consequences follow.
Comparing $\kappa$ values from collections with different relevance rates is
misleading. And $\kappa$ can be interpreted only beside the $p_o$ and $p_e$ it was
computed from.

\subsection*{Agreement on the ordering of systems}

The second form of the question sets the labels aside and asks about their use. Each
system in a set is scored twice with the same measure, once under each labeller's
grades, and the two resulting orderings of the systems are compared.

\emph{Kendall's $\tau$} compares them pair by pair~\cite{kendall1938tau}. For every
pair of systems, the two orderings either agree on which system is better, making the
pair \emph{concordant}, or disagree, making it \emph{discordant}. With $C$ concordant
and $D$ discordant pairs,
\begin{equation}
\tau = \frac{C - D}{C + D} .
\end{equation}
No two systems tie on the measure used here, so $C + D$ is simply the number of pairs,
and this coincides with Kendall's original coefficient. $\tau$ is 1 when the orderings
are identical, 0 when they are unrelated, and $-1$ when one reverses the other.

Agreement on labels and agreement on orderings are different quantities, and either
can be high while the other is low. Consider a labeller whose disagreements all run in
one direction, for example one that grades every non-relevant document one step too
high. It raises every system's score together, so it can leave the ordering nearly
intact while its $\kappa$ is low.

\subsection*{Resolving power}

A high $\tau$ says that most pairs of systems keep their order. It does not say which
pairs do not. Those are, predictably, the pairs whose scores are close. The ability of an evaluation to tell such close systems apart is its \emph{resolving power}, and Upadhyay et al.\ found UMBRELA's to be limited~\cite{upadhyay2025umbrela}.

Buckley and Voorhees give a way to measure it~\cite{buckley2000stability}. They
compared the orderings that different evaluations of the same systems produce, in
their case the same measure over different query formulations and over topic sets of
varying size. For each pair of systems they counted how often the two evaluations
disagreed about which was the better, and took the minority verdict as the error; the
\emph{error rate} is the total of those errors over all the decisions made, ties
included. Close pairs are set aside by a \emph{fuzziness value}: the percentage difference between two scores below which the scores are treated as equivalent. The closest pairs are the ones
most likely to be misordered, so raising the fuzziness value lowers the error rate, at
the price of leaving more pairs undecided. Together, a fuzziness value and the error
rate it yields state how large a difference must be before the ordering it implies can
be trusted.

Two questions therefore remain for the judge this thesis uses: whether a judge that
disagrees with the assessors on individual labels can still be trusted to rank
systems, and how small a difference it can resolve. Section~\ref{sec:met-judge-validation}
settles both.

\section{Generation evaluation metrics}\label{sec:bg-generation-metrics}

Retrieval measures say whether the right documents were found; they say nothing
about what is then done with them. In \emph{retrieval-augmented generation} a
language model --- the \emph{generator} --- receives the retrieved documents as
its \emph{context} alongside the question and writes the answer that is scored, so
the answer needs measures of its own. Two questions arise once it exists, and
neither answers the other: whether the answer is grounded in the context it was
given, and whether the answer is correct. An answer can be grounded in a context
that is itself wrong, and it can be correct from knowledge no context supplied.
This section defines a measure for both.

Faithfulness is defined by RAGAS, a framework that scores the stages of a
retrieval-augmented pipeline with a language model in place of a human
annotator~\cite{es2024ragas}; factual correctness was added later to its
implementation~\cite{ragaslib}, in the version named in
Section~\ref{sec:gen-method}. The two are built alike, so the construction is
given once and each measure follows from it. A text is first decomposed into
\emph{claims} --- standalone assertions, each of which can be checked on its own.
Every claim is then judged against some body of text, scoring 1 where it can be
inferred from that text and 0 where it cannot. The language model performs both
steps. What separates one measure from another is which text is decomposed and
what its claims are checked against.

\emph{Faithfulness} answers the first question: the generated answer is
decomposed, and its claims are checked against the retrieved context, taken as one
body of text. \emph{Factual correctness} answers the second, and needs something
to be correct about --- a \emph{reference}, a text holding the answer the question
should have received. The answer is decomposed again, and its claims are checked
against the reference instead. Its precision reading,
\emph{factual-correctness precision} (FC-P), is the share of those claims the
reference supports. The two measures are therefore one formula read twice,
\begin{equation}
\begin{aligned}
\mathrm{Faithfulness} &=
\frac{\bigl|\{c \in C(\text{answer}) : c \text{ inferable from the context}\}\bigr|}
     {\bigl|C(\text{answer})\bigr|} ,\\[6pt]
\text{FC-P} &=
\frac{\bigl|\{c \in C(\text{answer}) : c \text{ inferable from the reference}\}\bigr|}
     {\bigl|C(\text{answer})\bigr|} ,
\end{aligned}
\end{equation}
where $C(x)$ denotes the set of claims the text $x$ decomposes into. Both
denominators count the answer's own claims; only the text the claims are checked
against changes.

What each measure leaves out follows from that. Faithfulness requires no correct
answer to compare against. That is why it can be computed wherever a context
exists, and why it is not a quality score. It measures grounding rather than
truth: an answer scores 1 as long as the context supports every claim it makes,
even when the context is wrong. Factual-correctness precision is silent in the
other direction. It never sees the context, so an answer the generator wrote from
what it already knew scores as well as one the context supplied, provided the
reference supports it.

The denominator is where faithfulness is most easily misread, and the misreading
is worth naming. The fraction is taken over the claims of the \emph{answer}, never
over those of the context, so nothing about how much of the context was used
enters the score: an answer that leans on one line of a long context and an answer
that draws on all of it can both score 1. Faithfulness is a proportion of what was
said, not of what was available.

One property of the shared construction bears on every score it produces. How
finely a text is cut into claims is not fixed but governed by two parameters:
\emph{atomicity}, how far a sentence is split into separate assertions, and
\emph{coverage}, how much of a sentence's detail survives the split. Numerator and
denominator are both counts of claims, so the granularity of the decomposition
moves the whole fraction, and two answers of identical content can score
differently if they are cut differently. That is a property of the instrument, not
of the answers.

Faithfulness needs only a context. Factual correctness needs a reference, and for
a mathematics question the collection does not hand one over ready-made: the same
question usually admits several valid routes to the same result, so taking any one
correct answer as the reference would mark a correct response wrong for having
taken a different route. What a reference is built from instead is settled in
Section~\ref{sec:gen-method}.
